% =====================================================================
%  dati-gruppo.tex  --  Dati che valgono per TUTTI i documenti.
%  La classe gpt.cls li legge da sola: non dovete includere questo file.
%
%  ATTENZIONE: qui vanno solo dati che NON cambiano nel tempo.
%  Il capitolato e chi era presente a una riunione si scrivono nel
%  singolo documento: se li metteste qui, cambierebbero anche tutti i
%  PDF già approvati. (Se ci provate, la compilazione si ferma.)
% =====================================================================

\nomegruppo{Git Push Team}
\motto{--force \& pray}            % per ora non viene stampato
\numerogruppo{9}
\emailgruppo{gitpushteam@gmail.com}
\annoaccademico{2026/2027}

% Destinatari predefiniti. Un documento può cambiarli scrivendo nel suo
% preambolo, per esempio:  \destinatari{Prof.~Tullio Vardanega\newline Eggon S.r.l.}
\destinatari{Prof.~Tullio Vardanega\newline Prof.~Riccardo Cardin}  % \newline = a capo

% ---------------------------------------------------------------------
%  Persone del gruppo:   \membro{chiave}{Nome Cognome}{matricola}
%
%  - La "chiave" è un nome breve, tutto minuscolo, senza spazi né
%    accenti: è quella che scrivete nei documenti, es. \redattore{guglielmo}.
%  - NON cancellate mai una riga, nemmeno se qualcuno lascia il gruppo:
%    i verbali vecchi la usano ancora. Toglietelo invece da \componenti{...}
%    nei modelli dei verbali (template/modelli/).
%  - Esiste anche la chiave speciale "tutti", che stampa "Tutto il gruppo".
% ---------------------------------------------------------------------
\membro{biasini}{Enrico Biasini}{2147948}
\membro{cescobolla}{Luca Cesco Bolla}{2138006}
\membro{corro}{Alessio Corrò}{2145607}
\membro{guglielmo}{Francesco Guglielmo}{2137975}
\membro{mazzurana}{Francesco Mazzurana}{2149463}
\membro{spagnoli}{Gaetano Spagnoli}{2147977}
